\FloatBarrier

\section{Conclusions}

\begin{enumerate}

\item For LTCC, against the complete original hardware, the optimized CAD bundle contains 93.9\% fewer
triangles and 93.9\% fewer bytes. CAD-only loading CPU time decreases by 92.6\%. This is direct evidence
that the simplified files reduce loading work.

\item Complete LTCC initialization CPU time falls by 63.5\% and startup peak memory falls by 56.8\%. These
gains include both optimized files and the simplified native hierarchy.

\item Nonoptical electron transport is 1.29 times faster; hadron workloads take 1.35--1.53 times as much CPU
time. The separate cone-only measurements quantify the isolated Winston-cone substitution.

\item Full-port optical electron transport requires 1.73 times as much median CPU time in this benchmark. The
increased CPU use may reflect more complete transport after geometry defects are removed: the original model
repeatedly requires corrective navigation pushes, while optimized diagnostic samples contain more tracks and
steps. Lower original CPU time is therefore not sufficient evidence of faster correct transport. Revised
materials and containment also contribute, so the separate causes remain unresolved.

\item The optimized meshes remove 857 exact zero-area triangles and eliminate the 1,714 invalid/rejected-facet
warnings per original import. Navigation reports across the tested workloads fall from 9,546 to 1. Overlap and
containment reports remain, so these gains do not establish overlap-free geometry.

\item Ported TEXT and SQLite geometries agree for all six run-dependent configurations. Warning counts are
reported for the tested construction sample and transport workloads; they do not establish that the complete
CLAS12 geometry is free of overlaps. Detector-response, material-budget, magnetic-field, and full-experiment
comparisons remain outside this study.

\item The separate beamline apex pair contains 45.4\% fewer triangles while preserving topology, sharp hub
features, materials, and placements, with sampled shape differences below 0.05 mm. Its median loading wall
time falls by 15.7\% and startup peak memory by 17.4\%. Transport wall time improves by 1.77--2.01 times in
the first, smaller matched-primary sample. These gains describe the isolated GEMC3 pair and do not establish
full-experiment performance or detector-response equivalence.

\item The downstream tube pair is reduced by 3.75\% in triangles, and Segment 1 is closed by registering and
bridging its missing annular joins while retaining the bore. Import surface-hole warnings are removed, but
overlaps and navigation warnings remain. The measured transport ratios describe this limited workload and do
not establish response equivalence or faster correct transport.

\item The already-coarse adapter and torus exit-shield group permits only a 0.434\% triangle reduction with
protected casing mating faces. The adapter is renamed consistently, and the lead fill is retained exactly.
Its repeated benchmark does not establish substantial performance gains; shape preservation takes priority
over a larger triangle reduction.

\item The FT adjuster trio is replaced by one 48-triangle closed prism, reducing triangles by 98.74\% while
matching the estimated sealed reference union volume within 0.00003\%. The material and common placement are
retained, and internal overlaps disappear in the isolated check. The reference mass remains approximate
because two original shells are open; deliberate material redistribution prevents a claim of
detector-response equivalence.

\item The renamed FT inner shield and flange contains 36.14\% fewer triangles with retained bounds, alloy
material, placement, two closed components, and mounting/bore topology. Its small enclosed pocket is
unchanged, sampled surface differences remain below 0.05 mm, and total signed volume changes by about
-0.000044\%. The isolated benchmark quantifies loading, transport, and diagnostics without establishing
full-experiment response equivalence.

\item Nine renamed FT inner-shield bolts use smooth solid shafts and round heads, reducing triangles by 89.38\%
while retaining each copper volume within 0.0001\%. Shortened tips leave 0.15 mm gaps at the blind-hole
ends; independent collision, containment, and 10,000-point Geant4 checks find no shield interference.
Electron, pion, and neutron transport timings improve in this isolated sample, while proton tracking is
essentially unchanged. These results do not establish full-experiment speedup or response equivalence.

\item The renamed FT outer flange closes twelve bolt tunnels and reduces triangles by 60.79\% while retaining
the
central opening, engineering rim coordinates, bounds, and placement. The authorized hole closure adds
8.296635\% stainless-steel volume. Repeated tracking timings and residual assembly overlaps are recorded
separately; added material prevents attributing timing differences solely to fewer triangles or inferring
response equivalence.

The separate control at fixed material volume gives median tracking speedups of 1.073, 1.016, 1.034, and
1.013 for electrons, protons, positive pions, and neutrons, with the simplified mesh faster in every seed
pair of this larger sample.

\item The renamed FT outer mount repairs four triangular surface gaps and reduces triangles by 11.44\%,
retaining
its existing openings and placement. Protected radial relief removes about 0.24924\% material against the
sealed reference and provides at least 0.03597 mm radial separation from the flange in their common axial
interval. Surface checks and the 10,000-point assembly test find no joint overlap, resolving the earlier
flange/mount interference. Repeated tracking results show modest electron, proton, and pion gains with
slower neutron tracking; unrelated full-beamline diagnostics remain.

\item The final downstream FT inner shield has smooth flange faces with every bolt recess filled, while
retaining
both longitudinal grooves and the central passage. It contains 6,500 triangles, 85.27\% fewer than the raw
input and 51.85\% fewer than the first revision. Recess closure adds 0.436019\% tungsten volume relative to
the prior closed mesh. Upgrade and smooth-reference tracking comparisons are reported separately with clean
isolated transport; the subsequent Moller-cone revision resolves the shield/cone interference.

\item The renamed Moller cone fills its surface openings, retains its beam passage and outer bounds, and
reduces
triangles by 68.32\%. A widened protected counterbore removes about 0.168806\% tungsten relative to the
hole-filled reference and provides at least 0.187891 mm radial shield clearance. The cone/shield assembly
passes the 10,000-point overlap check. Raw-input timing includes solid and navigation warnings; the clean
clearance-reference control separately measures tessellation gains at effectively unchanged material volume.

\item Local relief of both Segment 1 sockets resolves the remaining downstream adapter and tube contacts
without
changing placements or neighboring meshes. Segment 1 has 5,096 triangles, 0.604764\% less steel than its
previous closed mesh, and at least 0.1 mm socket-face clearance. The three-part high-resolution overlap test
and full default beamline construction pass without exceptions. Separate repeated tracking results include
the repair and material change.

\item The beamline runtime database now selects every configuration by run number under \texttt{default}.
Run 16000 loads all 32 geometry volumes, and all 17 configured runs pass construction with unchanged
physical geometry. This selection fix has no independently measured tracking-speed gain.

\item The renamed Moller cone tip has 63.58\% fewer triangles with retained topology, sharp faces, bounds, and
a -0.000486\% material-volume change. Isolated tracking is 1.659, 1.133, 1.177, and 1.159 times faster for
electrons, protons, positive pions, and neutrons in this sample. Its pre-existing main-cone overlap remains in
the four affected complete configurations and is reported separately.

\item The renamed FT flange retains its opening, flat faces, bounds, and topology with 540 triangles, 2.17\%
fewer than the already-coarse original. Steel volume changes by -0.000109\%, and sampled surface differences
remain below 0.05 mm. Median tracking ratios remain within 0.18\% of one, providing no meaningful speed-gain
evidence. Existing neighboring assembly defects remain and are reported separately.

\item The renamed inner Moller-cone extension restores its missing radial-port wall and reduces triangles by
30.77\% while retaining its passages, bounds, and protected faces. Decimation changes the repaired reference
volume by +0.000009303\%, with sampled deviations below 0.016 mm. Repaired-reference tracking ratios are
1.261, 1.086, 1.144, 1.095 for electrons, protons, positive pions, and neutrons. The original is open and
produces navigation warnings; those raw timings are reported separately. Existing assembly overlaps remain.

\item The moller-cone extension lead fill retains its mating surface and openings with 24.49\% fewer
triangles. Median isolated tracking ratios are 1.0826, 1.0795, 1.1226, 1.0882 for electrons, protons, positive
pions, and neutrons. Every paired measurement improves in this sample. Existing assembly overlaps remain.

\item The moller-cone extension lead case retains its mating surface and openings with 42.40\% fewer
triangles. Median isolated tracking ratios are 1.0529, 1.0377, 1.0098, 1.0198 for electrons, protons, positive
pions, and neutrons. Some paired ranges include one; those particle samples do not establish a consistent
gain. Existing assembly overlaps remain.

\item The renamed Moller-cone extension tungsten add-on restores its missing radial-port wall and reduces
triangles by 56.80\% while retaining its passages, bounds, and protected faces. Decimation changes the
repaired reference volume by -0.0000046193\%, with sampled deviations below 0.013 mm. Repaired-reference
tracking ratios are 1.647, 1.548, 1.642, 1.535 for electrons, protons, positive pions, and neutrons. The
original is open and produces navigation warnings; those raw timings are reported separately. Existing
assembly overlaps remain.

\item Local FT/cone joint reliefs clear all eight checked interfaces while retaining the shield, flange, and
cone tip. Five meshes lose 0.119--1.209\% material by volume. All affected complete configurations construct
without exceptions, and final corrected tracking runs are clean. Isolated tracking ratios 0.8568, 0.8541,
0.8465, 0.8920 describe the eight-part clearance revision; they are not pure decimation gains.

\item The three HTCC cones have 49.20\% fewer triangles, and the entry cone's 144 boundary edges are closed.
All six TEXT/CAD and SQLite/CAD initialization checks and the three native/CAD database comparisons pass. The
Moller cone retains 24 non-manifold edges, and the cured entry cone introduces an observed contact of
approximately 0.053 mm with the entry window. These unresolved defects preclude a claim of complete geometry
validation; no performance or detector-response gain has been measured for this transfer.

\end{enumerate}
